% --------------------------------------------------------------- METHOD ----
\section{Method: EquiPhase DEQ}
\label{sec:method}

\paragraph{Learned potential with structural guarantees.}
EquiPhase DEQ represents a system by a scalar potential and reaches
equilibria by integrating damped conservative dynamics. We describe the
model in three deliberately separated layers, so that the reader can always
tell which properties are \emph{guaranteed by construction}, which are
\emph{inherited from an analytical prior}, and which are \emph{empirical}.
(i)~\textit{Structure:} the potential is a function of an embedding chosen
so that required invariances hold exactly — in the synthetic system, the
symmetry and identity properties are architectural (gates G1/G2$'$ below);
on dihedral data, the embedding
$[\sin\phi,\cos\phi,\sin\psi,\cos\psi]$ makes $V_\theta$ exactly
$2\pi$-periodic. (ii)~\textit{Prior:} where an analytical potential
$V_{\mathrm{base}}$ exists (synthetic system) we learn a residual,
$V_{\mathrm{total}}=V_{\mathrm{base}}+v_{\mathrm{net}}$; where none exists
(molecular data) $V_\theta$ is learned alone, and we disclose that the
prior layer is only partially applicable. (iii)~\textit{Empirics:} all
remaining claims are established by preregistered gates.

\paragraph{Damped velocity Verlet dynamics.}
States are phase-space pairs $(q,p)$. One step of the model applies
\begin{align}
q_{1/2} &= q_t + \tfrac{\Delta t}{2m}\,p_t, \\
p_{t+1} &= (1-\eta)\,p_t - \Delta t\,\nabla_q V(q_{1/2}),
\qquad \eta = \gamma \Delta t / m, \\
q_{t+1} &= q_{1/2} + \tfrac{\Delta t}{2m}\,p_{t+1}.
\end{align}
With $\eta=0$ this is the classical velocity Verlet integrator; with
$\eta\in(0,1)$ it contracts phase-space volume at a fixed, known rate while
retaining the integrator's geometric stability. Equilibria of the map are
exactly the pairs $(q^\star, 0)$ with $\nabla_q V(q^\star)=0$: fixed points
of the dynamics coincide with critical points of the learned potential, so
counting attractors of the model is counting metastable states of $V$.

% --------------------------------------------------------------- THEORY ----
\section{Why uniqueness guarantees exclude multistability}
\label{sec:theory}

\begin{proposition}[Dissipative structure]
\label{prop:dissipative}
Let $J$ be the Jacobian of one damped velocity Verlet step with damping
$\eta$, and $\Omega$ the canonical symplectic form. Then
$J^\top \Omega J = (1-\eta)\,\Omega$.
\end{proposition}
The step is therefore not symplectic (it does not conserve energy — by
design) but \emph{conformally} so: phase-space volume shrinks by the
constant factor $(1-\eta)$ per step, which is what drives trajectories into
the basins of $V$ rather than orbiting them. In the synthetic system this
identity is verified numerically to the reported gate tolerance rather than
assumed.

\begin{proposition}[Banach collapse]
\label{prop:banach}
If a DEQ update map $g$ is a contraction, $\mathrm{Lip}(g)\le c<1$, on a
complete metric space, then $g$ admits exactly one fixed point, and every
initialization converges to it.
\end{proposition}

\begin{corollary}
A DEQ whose well-posedness is certified via global contraction (or strong
monotonicity implying it) has attractor count $N=1$ independent of the data
it was trained on. Representing $N\ge2$ metastable states requires forgoing
the global uniqueness certificate.
\end{corollary}
EquiPhase forgoes exactly that certificate and replaces it with a weaker,
physically meaningful one: dissipativity (Proposition~\ref{prop:dissipative})
plus empirical convergence diagnostics reported per run. Section~\ref{sec:ala2}
shows both halves of this trade empirically: our verified-contraction
baseline collapses precisely as Proposition~\ref{prop:banach} mandates,
while EquiPhase recovers the multistable structure of real molecular data.

% ------------------------------------------------------- EXPERIMENT 1 ------
\section{Experiment 1: Synthetic Double-Well with Analytical Ground Truth}
\label{sec:toy}

We first validate on a 32-dimensional anisotropic double-well potential
whose minima ($\pm\sqrt{\alpha}\,e_1$), saddle points, and barrier height
($0.2275$) are known in closed form, so every empirical claim can be scored
against an analytical answer. The residual network $v_{\mathrm{net}}$ is a
small MLP ($34{\to}64{\to}1$; $2305$ parameters). Under a preregistered
gate battery (seed 7777): the structural symmetry and identity gates hold to
tolerances of order $10^{-9}$; two attractors are recovered
($N=2$) with contraction factors $\rho = 0.9626/0.9628$ against the
corrected theoretical value $0.96273$; and the measured escape barrier
$0.230111$ exceeds the analytical barrier $0.2275$, consistent with the
learned residual steepening the well walls. Two preregistered stress
disclosures accompany these results: at the boundary setting $\alpha=1.2$
one gate exceeds its per-$\alpha$ limit by $1.0\times10^{-5}$, and $1$ of
$100$ novel-seed initializations diverges; both are reported as limits of
the method rather than removed. Under the same protocol, an unconstrained
vanilla DEQ fails to converge ($0/100$ initializations; residual
$\approx5.9\times10^{-3}$, $\rho\approx0.9999$), and a monotone DEQ
collapses to a single attractor with its training loss frozen — the same
qualitative signature that reappears on molecular data in
Section~\ref{sec:ala2}.
